% GENERATED FILE. Edit canonical lessons, then run npm run generate:books.
% canonical-source-hash: fnv1a64:9f0dd8ebfef2fd35
% canonical-lessons: KA-C339-galabhe, KA-C339-kolahala, KA-C339-janasandani, KA-C339-caduru, KA-C339-kaltulita

\chapter{A Riot in the Market}
\label{ch:ka-a-riot-in-the-market}
\hlchaptermodality{339}

\begin{chapteropening}
\textbf{By the end of this chapter:} \emph{I can say a riot, an uproar, a crowd, to scatter and a stampede.}

Complete the last lesson of chapter 339: I can say a riot, an uproar, a crowd, to scatter and a stampede.
\end{chapteropening}

\section[\texorpdfstring{\kn{ಗಲಭೆ} (galabhe)}{ಗಲಭೆ (galabhe)}]{\kn{ಗಲಭೆ} (galabhe) — a riot; a disturbance}

\label{lesson:KA-C339-galabhe}



\begin{quote}
Before the new one: say the Kannada for the truth, then the Kannada for to listen closely.
\end{quote}

\subsection*{You'll want to know: \kn{ಗಲಭೆ}}
\textbf{\kn{ಗಲಭೆ}} — \emph{galabhe} — \textquotedblleft{}a riot; a disturbance\textquotedblright{}.

\begin{quote}
\kn{ಸಂತೆಯಲ್ಲಿ} \kn{ಸಣ್ಣ} \kn{ಜಗಳ} \kn{ದೊಡ್ಡ} \kn{ಗಲಭೆ} \kn{ಆಯಿತು}.
\end{quote}

\begin{quote}
\emph{santeyalli saṇṇa jagaḷa doḍḍa galabhe āyitu.}
\end{quote}

\begin{quote}
A small quarrel at the market became a big riot.
\end{quote}

Plural: \kn{ಗಲಭೆಗಳು}, \emph{galabhegaḷu}.

In the news: \kn{ಗಲಭೆ} \kn{ನಡೆಯಿತು}, there was a riot.

\begin{cousinweb}
\textbf{\kn{ಗಲಭೆ}} (\emph{galabhe}) is native Dravidian: DEDR 1310 lists Kannada \emph{galabe}, \textquotedblleft{}hubbub, clamour\textquotedblright{}, and \emph{galabili}, \textquotedblleft{}disorder\textquotedblright{}, beside Telugu \emph{galaba}, \textquotedblleft{}confusion, noise\textquotedblright{}.
\end{cousinweb}

\begin{grammarlens}[title={What you've built}]
One more word for telling what happened.
\end{grammarlens}

\subsection*{Your turn}
\begin{itemize}
\raggedright
  \item \emph{Say it:} \emph{galabhe}
  \item \emph{Say it:} \emph{galabhe}, once more
  \item \emph{Say it:} \emph{santeyalli saṇṇa jagaḷa doḍḍa galabhe āyitu.} — A small quarrel at the market became a big riot.
  \item \emph{Recall:} say the Kannada for the truth, then the Kannada for to listen closely, then say \emph{galabhe} again
\end{itemize}

\begin{tcolorbox}[breakable,colback=teal!4,colframe=teal!35!black,title={Before you move on}]
What does \kn{ಗಲಭೆ} mean? (\textquotedblleft{}A riot; a disturbance\textquotedblright{}.) Say it once more.
\end{tcolorbox}
\section[\texorpdfstring{\kn{ಕೋಲಾಹಲ} (kōlāhala)}{ಕೋಲಾಹಲ (kōlāhala)}]{\kn{ಕೋಲಾಹಲ} (kōlāhala) — an uproar; a commotion}

\label{lesson:KA-C339-kolahala}



\begin{quote}
Before the new one: say the Kannada for to listen closely, then the Kannada for a riot.
\end{quote}

\subsection*{You'll want to know: \kn{ಕೋಲಾಹಲ}}
\textbf{\kn{ಕೋಲಾಹಲ}} — \emph{kōlāhala} — \textquotedblleft{}an uproar; a commotion\textquotedblright{}.

\begin{quote}
\kn{ಗೆದ್ದ} \kn{ಸುದ್ದಿ} \kn{ಕೇಳಿ} \kn{ಮೈದಾನದಲ್ಲಿ} \kn{ಕೋಲಾಹಲ} \kn{ಎದ್ದಿತು}.
\end{quote}

\begin{quote}
\emph{gedda suddi kēḷi maidānadalli kōlāhala edditu.}
\end{quote}

\begin{quote}
Hearing of the win, an uproar rose in the ground.
\end{quote}

An uproar \textquotedblleft{}rises\textquotedblright{}: \kn{ಕೋಲಾಹಲ} \kn{ಏಳು}. It can be joyful or angry.

\begin{cousinweb}
\textbf{\kn{ಕೋಲಾಹಲ}} (\emph{kōlāhala}) is borrowed from Sanskrit \emph{kolāhala}, \textquotedblleft{}a loud and confused sound, uproar\textquotedblright{} (Monier-Williams), a word that imitates its own noise.
\end{cousinweb}

\begin{grammarlens}[title={What you've built}]
One more word for telling what happened.
\end{grammarlens}

\subsection*{Your turn}
\begin{itemize}
\raggedright
  \item \emph{Say it:} \emph{kōlāhala}
  \item \emph{Say it:} \emph{kōlāhala}, once more
  \item \emph{Say it:} \emph{gedda suddi kēḷi maidānadalli kōlāhala edditu.} — Hearing of the win, an uproar rose in the ground.
  \item \emph{Recall:} say the Kannada for to listen closely, then the Kannada for a riot, then say \emph{kōlāhala} again
\end{itemize}

\begin{tcolorbox}[breakable,colback=teal!4,colframe=teal!35!black,title={Before you move on}]
What does \kn{ಕೋಲಾಹಲ} mean? (\textquotedblleft{}An uproar; a commotion\textquotedblright{}.) Say it once more.
\end{tcolorbox}
\section[\texorpdfstring{\kn{ಜನಸಂದಣಿ} (janasandaṇi)}{ಜನಸಂದಣಿ (janasandaṇi)}]{\kn{ಜನಸಂದಣಿ} (janasandaṇi) — a crowd; a throng}

\label{lesson:KA-C339-janasandani}



\begin{quote}
Before the new one: say the Kannada for a riot, then the Kannada for an uproar.
\end{quote}

\subsection*{You'll want to know: \kn{ಜನಸಂದಣಿ}}
\textbf{\kn{ಜನಸಂದಣಿ}} — \emph{janasandaṇi} — \textquotedblleft{}a crowd; a throng\textquotedblright{}.

\begin{quote}
\kn{ಹಬ್ಬದ} \kn{ದಿನ} \kn{ದೇವಸ್ಥಾನದ} \kn{ಮುಂದೆ} \kn{ಜನಸಂದಣಿ} \kn{ಇತ್ತು}.
\end{quote}

\begin{quote}
\emph{habbada dina dēvasthānada munde janasandaṇi ittu.}
\end{quote}

\begin{quote}
On the festival day there was a crowd in front of the temple.
\end{quote}

\kn{ಜನ}, people, and \kn{ಸಂದಣಿ}, crowding: a press of people.

\begin{cousinweb}
\textbf{\kn{ಜನಸಂದಣಿ}} (\emph{janasandaṇi}) is a Kannada compound of \kn{ಜನ}, \textquotedblleft{}people\textquotedblright{}, from Sanskrit \emph{jana} (Monier-Williams), and \kn{ಸಂದಣಿ}, \textquotedblleft{}crowding\textquotedblright{}; the origin of \kn{ಸಂದಣಿ} is not settled.
\end{cousinweb}

\begin{grammarlens}[title={What you've built}]
One more word for telling what happened.
\end{grammarlens}

\subsection*{Your turn}
\begin{itemize}
\raggedright
  \item \emph{Say it:} \emph{janasandaṇi}
  \item \emph{Say it:} \emph{janasandaṇi}, once more
  \item \emph{Say it:} \emph{habbada dina dēvasthānada munde janasandaṇi ittu.} — On the festival day there was a crowd in front of the temple.
  \item \emph{Recall:} say the Kannada for a riot, then the Kannada for an uproar, then say \emph{janasandaṇi} again
\end{itemize}

\begin{tcolorbox}[breakable,colback=teal!4,colframe=teal!35!black,title={Before you move on}]
What does \kn{ಜನಸಂದಣಿ} mean? (\textquotedblleft{}A crowd; a throng\textquotedblright{}.) Say it once more.
\end{tcolorbox}
\section[\texorpdfstring{\kn{ಚದುರು} (caduru)}{ಚದುರು (caduru)}]{\kn{ಚದುರು} (caduru) — to scatter; to disperse}

\label{lesson:KA-C339-caduru}



\begin{quote}
Before the new one: say the Kannada for an uproar, then the Kannada for a crowd.
\end{quote}

\subsection*{You'll want to know: \kn{ಚದುರು}}
\textbf{\kn{ಚದುರು}} — \emph{caduru} — \textquotedblleft{}to scatter; to disperse\textquotedblright{}.

\begin{quote}
\kn{ಪೊಲೀಸರು} \kn{ಬಂದ} \kn{ಕೂಡಲೇ} \kn{ಜನ} \kn{ಚದುರಿದರು}.
\end{quote}

\begin{quote}
\emph{polīsaru banda kūḍalē jana caduridaru.}
\end{quote}

\begin{quote}
As soon as the police came, the people scattered.
\end{quote}

Past (it …): \kn{ಚದುರಿತು}, \emph{caduritu}.

A crowd or clouds disperse: \kn{ಮೋಡ} \kn{ಚದುರಿತು}. To scatter them is \kn{ಚದುರಿಸು}.

\begin{cousinweb}
\textbf{\kn{ಚದುರು}} (\emph{caduru}) is native Dravidian: DEDR 1546 lists Kannada \emph{cadaṟu}, \textquotedblleft{}to be scattered about, disperse\textquotedblright{}, beside Tamil \emph{citar}, \textquotedblleft{}to scatter\textquotedblright{}.
\end{cousinweb}

\begin{grammarlens}[title={What you've built}]
One more word for telling what happened.
\end{grammarlens}

\subsection*{Your turn}
\begin{itemize}
\raggedright
  \item \emph{Say it:} \emph{caduru}
  \item \emph{Say it:} \emph{caduru}, once more
  \item \emph{Say it:} \emph{polīsaru banda kūḍalē jana caduridaru.} — As soon as the police came, the people scattered.
  \item \emph{Recall:} say the Kannada for an uproar, then the Kannada for a crowd, then say \emph{caduru} again
\end{itemize}

\begin{tcolorbox}[breakable,colback=teal!4,colframe=teal!35!black,title={Before you move on}]
What does \kn{ಚದುರು} mean? (\textquotedblleft{}To scatter; to disperse\textquotedblright{}.) Say it once more.
\end{tcolorbox}
\section[\texorpdfstring{\kn{ಕಾಲ್ತುಳಿತ} (kāltuḷita)}{ಕಾಲ್ತುಳಿತ (kāltuḷita)}]{\kn{ಕಾಲ್ತುಳಿತ} (kāltuḷita) — a stampede}

\label{lesson:KA-C339-kaltulita}



\begin{quote}
Before the new one: say the Kannada for a crowd, then the Kannada for to scatter.
\end{quote}

\subsection*{You'll want to know: \kn{ಕಾಲ್ತುಳಿತ}}
\textbf{\kn{ಕಾಲ್ತುಳಿತ}} — \emph{kāltuḷita} — \textquotedblleft{}a stampede\textquotedblright{}.

\begin{quote}
\kn{ಜಾತ್ರೆಯಲ್ಲಿ} \kn{ಕಾಲ್ತುಳಿತ} \kn{ಆಗಿ} \kn{ಹಲವರಿಗೆ} \kn{ಗಾಯವಾಯಿತು}.
\end{quote}

\begin{quote}
\emph{jātreyalli kāltuḷita āgi halavarige gāyavāyitu.}
\end{quote}

\begin{quote}
There was a stampede at the fair and many were hurt.
\end{quote}

\kn{ಕಾಲು}, foot, and \kn{ತುಳಿತ}, trampling.

\begin{cousinweb}
\textbf{\kn{ಕಾಲ್ತುಳಿತ}} (\emph{kāltuḷita}) is a native Dravidian compound: \emph{kāl}, \textquotedblleft{}foot, leg\textquotedblright{} (DEDR 1479), and a noun from \emph{tuḷi}, \textquotedblleft{}to tread, to trample\textquotedblright{} (DEDR 3522).
\end{cousinweb}

\begin{grammarlens}[title={What you've built}]
One more word for telling what happened.
\end{grammarlens}

\subsection*{Your turn}
\begin{itemize}
\raggedright
  \item \emph{Say it:} \emph{kāltuḷita}
  \item \emph{Say it:} \emph{kāltuḷita}, once more
  \item \emph{Say it:} \emph{jātreyalli kāltuḷita āgi halavarige gāyavāyitu.} — There was a stampede at the fair and many were hurt.
  \item \emph{Recall:} say the Kannada for a crowd, then the Kannada for to scatter, then say \emph{kāltuḷita} again
\end{itemize}

\begin{tcolorbox}[breakable,colback=teal!4,colframe=teal!35!black,title={Before you move on}]
What does \kn{ಕಾಲ್ತುಳಿತ} mean? (\textquotedblleft{}A stampede\textquotedblright{}.) Say it once more.
\end{tcolorbox}
